\chapter{State and Input: Branches, Bitmasks, and Boundaries}
\label{ch:asm2}

In Chapter~\ref{ch:asm1} a counter advanced on its own. A game, though, responds
to a player. This chapter teaches the assembly you need to read the controller
and to make decisions: the input bitmask, the bitwise \code{andi} test, and the
conditional branch. By the end you will move an object with the joystick and keep
it on the screen: the core interactive loop of every game in this book. We follow
the structure of the platform's ASCII Pong update routine, which was written to be
read by first-year students \citep{prg32tutorialascii,prg32repo}.

\section{Decisions: the conditional branch}

So far our instructions ran in a straight line. To respond to the player, a
program must choose: if a button is pressed, move; otherwise, do not. In \riscv
this choice is a \emph{conditional branch}, an instruction that jumps to a label
only when a condition holds.

\begin{prgconcept}
A \textbf{conditional branch} compares registers and jumps to a label only if the
comparison is true; otherwise execution simply continues to the next instruction.
The branch is how a processor makes every decision, from ``is this button
pressed?'' to ``has the ball hit the wall?''
\end{prgconcept}

The branches you will use most are few. \code{beqz t, label} jumps if \reg{t} is
zero; \code{bltz t, label} jumps if \reg{t} is negative; \code{ble a, b, label}
jumps if \reg{a} $\le$ \reg{b}; and the unconditional \code{j label} always jumps.
A common idiom, which surprises beginners, is to branch on the \emph{opposite} of
what you want: to do something ``if the button is pressed,'' you test the button
bit and \code{beqz}, branch away, when it is \emph{not} pressed, letting the
movement code run only when it is. This reads backwards at first but becomes
natural quickly.

\begin{prgaside}
Local labels written as a bare number followed by a colon, like \code{1:}, and
referenced as \code{1f} (the next \code{1} ``forward'') or \code{1b} (``backward''),
are an assembler convenience for short jumps that need no memorable name. The
platform's examples use them heavily for the small skips inside an update routine
\citep{prg32repo}.
\end{prgaside}

\section{Reading the controller}

Recall from Chapter~\ref{ch:platform} that \prg returns all buttons at once as a
bitmask. In assembly you obtain it by calling \code{prg32\_input\_read}; the
result arrives, by convention, in \reg{a0}. The very first thing to do is copy it
somewhere safe, because the next \code{call} you make would overwrite \reg{a0}.

\lstinputlisting[
    style=asm,
    caption={%
        \fname{get\_input.S} \protect\\
        Read the input bitmask and keep it. The copy to a temporary is essential, because later calls reuse \reg{a0}.%    
    }
]{examples/asm_state_and_input/get_input.S}

\begin{prgwarn}
Copy the input mask out of \reg{a0} immediately. \reg{a0} is the argument and
return register; the moment you call another function, its old contents are gone.
The platform's examples copy the mask into a temporary on the very next line for
exactly this reason \citep{prg32tutorialascii}.
\end{prgwarn}

To test a single button you isolate its bit with \code{andi}, the
``and-immediate'' instruction, using the bit value from
Table~\ref{tab:buttons}. The result is non-zero exactly when that button is
pressed.

\lstinputlisting[
    style=asm,
    caption={%
        \fname{button\_test.S} \protect\\
        Testing one button. \code{andi} keeps only the chosen bit; \code{beqz} skips the body when that bit is clear.%    
    }
]{examples/asm_state_and_input/button_test.S}

\begin{prgconcept}
\code{andi t3, t0, 1} computes the bitwise AND of the input mask with the
constant \code{1}, leaving in \reg{t3} the value of bit~0 alone. If LEFT is
pressed, \reg{t3} is non-zero; if not, it is zero. Testing a button is therefore
``mask the bit, then branch on whether the result is zero.''
\end{prgconcept}

\section{Moving an object with the joystick}

We now assemble a complete \code{update} that moves an object left and right with
the joystick. The state is a single word in \code{.data}, exactly as in
Chapter~\ref{ch:asm1}. The routine reads input, tests LEFT and RIGHT, and updates
the position. For now, we limit the draw to displaying the position on screen.
    In the future, the position will make a graphics object move 
    
    \citep{prg32repo}.

    \lstinputlisting[
        style=asm,
        caption={%
            A full update: read input, move on LEFT/RIGHT, store the result. 
            Compare with the platform's \texttt{pong\_ascii\_update}.%
        }
    ]{examples/asm_state_and_input/input_read.S}

Trace it once by hand. Suppose \code{game\_x} holds \code{20} and the player
holds RIGHT. The input read returns a mask with bit~1 set. The LEFT test
(\code{andi ... 1}) yields zero, so \code{beqz} skips the left move. The RIGHT
test (\code{andi ... 2}) yields non-zero, so the right move runs and \reg{t2}
becomes \code{21}. The store writes \code{21} back. Next frame it becomes
\code{22}, and so on. The object slides right as long as RIGHT is held. This kind
of by-hand trace is the most valuable debugging skill you will build in this
course.

The code also introduces a tiny performance trick. Instead of redrawing the
screen on every update, it stores a small flag in \code{text\_update} whenever
LEFT or RIGHT actually changes \code{game\_x}. The draw routine checks this
flag and skips the clear/write sequence when nothing moved. In the previous
chapter, the example redrew continuously, but a real game should avoid doing
work when the state does not change.

There is also something to notice about the numbers on screen. If you move
left past zero, the value underflows and the display shows \code{0xffffffff} instead
of a normal count. If you move right past \code{9}, the console printer keeps
counting in hexadecimal, so \code{10} appears as \code{a}, \code{11} as
\code{b}, and so on. That is because we are displaying the 
32 bit score word is displayed as an hexadecimal value.

\section{Staying on the screen: the boundary clamp}

Left unchecked, our object will slide right off the edge of the screen and
vanish. The fix is a \emph{clamp}: after computing the new position, force it back
into the legal range. For the 40-column ASCII console the legal range is 0 to 39;
for the 320-pixel graphics viewport it will be 0 to a value that accounts for the
object's width. The logic is identical either way.

\lstinputlisting[
    style=asm,
    caption={%
        \fname{x\_clamping.S} \protect\\
        Clamping x to the range [0, 39]. Negative values snap to 0; values past the edge snap to the maximum.% %
    }
]{examples/asm_state_and_input/x_clamping.S}

\begin{prgconcept}
A \textbf{boundary clamp} keeps a value inside a legal range: if it went below the
minimum, set it to the minimum; if above the maximum, set it to the maximum;
otherwise leave it alone. Almost every off-screen bug in a beginner's game is a
missing or wrong clamp. When an object disappears, suspect the clamp first.
\end{prgconcept}

\begin{prgwarn}
The maximum is not the screen width; it is the screen width minus the object's
own size. A 64-pixel-wide paddle on a 320-pixel screen may go no further right
than $320-64 = 256$, or its right edge would leave the viewport. The platform's
graphics Pong clamps the paddle to exactly \code{256} for this reason
\citep{prg32repo}.
\end{prgwarn}

\begin{prgtry}
Add the clamp above to your \code{update}. Write
one sentence explaining which comparison protects the right edge. This
break-and-fix exercise is taken directly from the platform's labs
\citep{prg32tutorialgraphic}.
\end{prgtry}

\section{Debugging like a systems programmer}

When an assembly game misbehaves, the platform's guidance is explicit and worth
adopting as a discipline: do not guess, trace state \citep{prg32tutorial}. Ask, in
order, which register held the last correct value, which instruction changed it,
whether a \code{call} overwrote a temporary you expected to keep, and whether a
branch skipped the code that updates memory. Reduce the program to the smallest
version that still shows the bug. Nine times in ten the answer is one of three
things: a missing \reg{ra} save, an address-versus-value confusion, or a branch
that goes the wrong way.

\begin{prgcheck}
You should now be able to read the input bitmask into a register; test an
individual button with \code{andi} and \code{beqz}; move a value in \code{.data}
in response to the joystick; and clamp that value to a legal range. With these,
you have written the interactive core of a game. Chapter~\ref{ch:asm3} adds the
one missing ingredient: drawing it on the screen.
\end{prgcheck}
